% Generated draft: requires mathematical and source review.
\documentclass[11pt]{article}
\usepackage[margin=25mm]{geometry}
\usepackage{fontspec}
\setmainfont{Latin Modern Roman}
\usepackage{amsmath,amssymb,mathtools,unicode-math}
\setmathfont{Latin Modern Math}
\usepackage{enumitem,needspace,xurl,hyperref}
\setlength{\parindent}{0pt}
\setlength{\parskip}{5pt}
\setlength{\emergencystretch}{3em}
\newcommand{\N}{\mathbb{N}}
\newcommand{\Z}{\mathbb{Z}}
\newcommand{\Q}{\mathbb{Q}}
\newcommand{\R}{\mathbb{R}}
\newcommand{\C}{\mathbb{C}}
\newcommand{\F}{\mathbb{F}}
\newcommand{\A}{\mathbb{A}}
\newcommand{\PP}{\mathbb{P}}
\newcommand{\E}{\mathbb{E}}
\newcommand{\eps}{\varepsilon}
\newcommand{\dd}{\,\mathrm{d}}
\begin{document}
\section*{UnsolvedMath ID: 700733}
\textbf{Source problem code:} catalogue entry
\subsection{Definitions and mathematical statement}
A group \(G\) is torsion-free if \(g^n=1\), with \(g\in G\) and an integer \(n\geq 1\), implies \(g=1\). It has the unique product property, denoted \(\mathcal{UP}\), if for every pair of nonempty finite subsets \(A,B\subseteq G\), there exists \(c\in AB\) such that
\[
\bigl|\{(a,b)\in A\times B:ab=c\}\bigr|=1.
\]
For clarity, the classical presentation interpretation of the small cancellation condition is as follows. Let \(G\cong\langle X\mid R\rangle\), where \(R\) is a symmetrized set of nonempty cyclically reduced words in the free group on \(X\): it is closed under inverses and cyclic permutations. A piece is a nonempty word that is an initial segment of two distinct words of \(R\). Writing \(|w|\) for word length, the presentation satisfies \(C'(\lambda)\) if every piece \(p\) that is an initial segment of \(r\in R\) satisfies
\[
|p|<\lambda |r|.
\]
No finiteness assumption on \(X\) or \(R\) is supplied.

The source asks whether every torsion-free group satisfying \(C'(\lambda)\) for \(\lambda\ll 1\) has \(\mathcal{UP}\). Under the interpretation that this means a uniform sufficiently small positive parameter and the existence of a classical presentation as above, the question becomes: prove or disprove
\[
\begin{gathered}
\exists\lambda_0>0\ \ \forall\lambda\in(0,\lambda_0]\ \ \forall\langle X\mid R\rangle,\\
\left[\begin{array}{c}
G=\langle X\mid R\rangle\text{ is torsion-free},\\
\langle X\mid R\rangle\text{ satisfies }C'(\lambda)
\end{array}\right]
\ \Longrightarrow\ G\text{ has }\mathcal{UP}.
\end{gathered}
\]
This is an explicit interpretation of the source wording, requiring review: the record does not specify a numerical threshold, explicitly quantify \(\lambda\ll1\), or identify the intended small cancellation convention.

The source also states the consequent assertion that \(KG\) has no zero divisors. If \(K\) denotes a field, \(KG\) is the group algebra of finite formal sums \(\sum_{g\in G}a_g g\), with multiplication extending the group multiplication bilinearly; having no zero divisors means
\[
u,v\in KG,\quad uv=0\quad\Longrightarrow\quad u=0\ \text{or}\ v=0.
\]
The record does not specify the coefficient object \(K\); the field interpretation is conditional and must be checked against the source.
\subsection{Short English statement}
The question asks whether every torsion-free group with sufficiently strong small cancellation has the unique product property: any two nonempty finite subsets must have a product admitting exactly one representation. The source adds a consequence concerning zero divisors in \(KG\). The intended small cancellation convention, the quantification of \(\lambda\ll1\), and the coefficient object \(K\) require clarification.
\subsection{Sources}
\begin{itemize}
\item Supplied local UnsolvedMath registry record, permanent ID \(700733\), exact source problem code \texttt{catalogue entry}, title ``Kourovka Notebook Problem 11.40''; upstream revision \texttt{d371ee917ac9647fbb06cd6fc529deec6f3dda82}.
\item Attributed in the supplied record to S. V. Ivanov, proposed in 1990; original problem number 11.40 in \emph{The Kourovka Notebook (Unsolved Problems in Group Theory)}, 11th issue (1990). The record also identifies the 21st issue edition (2026), \texttt{arXiv:1401.0300v48}.
\item The supplied collection description names E. I. Khukhro and V. D. Mazurov as editors of the Kourovka Notebook collection covering Issues 1--21 (1965--2026).
\item The registry labels the problem ``open'' and records an update timestamp of \texttt{2026-10-16T14:00:00Z}. These are supplied registry metadata, not an independently verified assessment of the literature. The cited source has not been independently consulted for this draft.
\end{itemize}
\end{document}
